\subsection{以 B × I 为底空间的主纤维丛}

设 G 是一个拓扑群。我们将看到，在每个陈述中将“局部平凡纤维空间”替换为“群 G 的主纤维空间”后，命题 2 及其推论仍然有效。

引理 4. --- 设 B 是一个拓扑空间，G 是一个拓扑群，且 E、E′ 是以 B 为底、群 G 的主纤维丛。记 F 为拓扑空间 G，并赋予它由 G×G 给出的左作用 \((g,g')\cdot f=g'fg^{-1}\)，并记 \(M = (\mathrm{E}\times_{\mathrm{B}}\mathrm{E}')\times^{\mathrm{G}\times\mathrm{G}}\mathrm{F}\) 为与之相关联的、纤维型为 F 且底空间为 B 的局部平凡纤维丛。

层 \(\mathcal{S}\)（即 B 上由 M 的截面组成的层）同构于 B 上由 E 到 E' 的主纤维丛同构组成的层。

记 \(p\)、\(p'\) 和 \(q\) 为 B-空间 E、E' 和 M 的投影。对于 \((x, x') \in \mathrm{E} \times_{\mathrm{B}} \mathrm{E}'\) 和 \(f \in \mathrm{F}\)，记 \([x, x', f]\) 为元素 \(((x, x'), f) \in (\mathrm{E} \times_{\mathrm{B}} \mathrm{E}') \times \mathrm{F}\) 在 M 中的等价类。记 \(\mathcal{M}\) 为 B 上由 E 到 E' 的主纤维丛同构组成的层；它在 B 的开子集 U 上的截面是从 \(E_{U}\) 到 \(E_{U}'\) 的主纤维丛同构。

记 \(e\) 为 G 的单位元。设 U 是 B 的开子集，且 \(\varphi\colon\mathrm{E}_{\mathrm{U}}\to\mathrm{E}_{\mathrm{U}}^{\prime}\) 是群 G、底空间 U 的主纤维丛同构。从 \(\mathrm{E}_{\mathrm{U}}\) 到 \((\mathrm{E}\times_{\mathrm{B}}\mathrm{E}^{\prime})\times^{\mathrm{G}\times\mathrm{G}}\mathrm{F}\) 的映射由 \(x\mapsto[x,\varphi(x),e]\) 定义，且是连续的。对于 \(g\in\mathrm{G}\) 和 \(x\in\mathrm{E}_{\mathrm{U}}\)，它将 \(x\cdot g\) 送到 \([x\cdot g,\varphi(x\cdot g),e]=[x,\varphi(x),geg^{-1}]=[x,\varphi(x),e]\)。因而存在唯一连续映射 \(\alpha_{\mathrm{U}}(\varphi)\colon\mathrm{U}\to\mathrm{M}\)，使得 \(\alpha_{\mathrm{U}}(\varphi)(p(x))=[x,\varphi(x),e]\)，对于所有 \(x\in\mathrm{E}_{\mathrm{U}}\) 均成立；我们有 \(\alpha_{\mathrm{U}}(\varphi)\in\mathcal{S}(\mathrm{U})\)。

显然，映射 \(\alpha_{U}\) 定义了从层 M 到层 S 的层态射 \(\alpha\)。

引理 5. --- 层态射 α 是一个同构。

首先假设主纤维丛 E 和 E' 都是可平凡化的；取截面 \(i: B \to E\) 和 \(i': B \to E'\)。于是存在唯一连续映射 \(\theta\)，从 M 到 \(B \times G\)，使得

\[
\theta ([ i (b) \cdot g, i ^ {\prime} (b) \cdot g ^ {\prime}, f ]) = (b, g ^ {\prime} f g ^ {- 1})
\]

对 \(b \in \mathrm{B}\)、\(g \in \mathrm{G}\)、\(g' \in \mathrm{G}\) 和 \(f \in \mathrm{F}\)，该映射是连续的；它是 B-空间的同构。设 \(\varphi\) 是从 E 到 \(\mathrm{E}'\) 的主纤维丛同构；存在唯一连续映射 \(\gamma: \mathrm{B} \to \mathrm{G}\)，使得 \(\varphi(i(b)) = i'(b) \cdot \gamma(b)\) 对于 \(b \in \mathrm{B}\) 成立。\(\varphi\) 在映射 \(\alpha_{\mathrm{B}}\) 下的像是从 B 到 M 的映射 \(b \mapsto \theta^{-1}(b, \gamma(b))\)。这表明 \(\alpha_{\mathrm{B}}\) 是双射。

因此，\(\alpha_{\mathrm{U}}\) 对于 B 的每个开集 U 都是双射，只要主纤维丛 \(\mathrm{E_U}\) 和 \(\mathrm{E_U}'\) 是可平凡化的。由 I, p. 55 的 corollary 2，这意味着 \(\alpha\) 是层的同构。

命题 3. --- 设 G 是一个拓扑群，B 是一个仿紧拓扑空间，且 (E, p) 是以 B × I 为底、群 G 的主纤维丛。置 \(E_{1} = \overline{p}^{-1}(B \times \{1\})\)，并令 \(p_{1}: E_{1} \to B\) 为映射 \(pr_{1} \circ p \mid E_{1}\)。则 (E, p) 与 \((E_{1} \times I, p_{1} \times Id_{I})\) 是以 B × I 为底、群 G 的同构主纤维丛。

设 F 为拓扑空间 G，并赋予它由 \(G \times G\) 给出的左作用 \((g, g') \cdot f = g' fg^{-1}\)。设 \((M, q)\) 为以 \(B \times I\) 为底、与主纤维丛 \(E \times_{B \times I}(E_1 \times I)\) 相关联的、群为 \(G \times G\)、纤维型为 F 的局部平凡纤维丛。置 \(M_1 = \overline{q}^{-1}(B \times \{1\})\)

且 \(q_1 = \mathrm{pr}_1 \circ q \mid M_1\)；B-空间 \((M_1, q_1)\) 被识别为与 \(E_1 \times_\mathrm{B} E_1\) 相关联的、纤维型为 F 的局部平凡纤维丛。由引理 4（其中取主纤维丛 E 和 \(E'\) 都等于 \(E_1\)），B-空间 \(M_1\) 有一个截面。由于 \(B \times I\)-空间 \((M, q)\) 与 \((M_1 \times I, q_1 \times \mathrm{Id}_{I})\) 同构（IV, p. 440, Prop. 2），\(B \times I\)-空间 \((M, q)\) 有一个截面，这意味着主 G-丛 E 和 \(E_1 \times I\) 同构。

推论 1. --- 设 B 是一个仿紧拓扑空间，G 是一个拓扑群，且 (E,p) 是 B × I 上的主 G-丛。对于 \(t \in I\)，记 \((E_{t}, p_{t})\) 为主 B-纤维丛 \(i_{t}^{*}E\)，其中 \(i_{t}: B \to B \times \{t\}\) 是由 \(b \mapsto (b, t)\) 给出的映射。主 B-纤维丛 \(E_{0}\) 和 \(E_{t}\) 对每个 \(t \in I\) 都同构。

推论 2. --- 设 B 和 B' 是拓扑空间，G 是一个拓扑群，E 是群 G、底空间 B 的主 B-纤维丛。令 \(f_0\) 和 \(f_1\) 为从 B' 到 B 的连续映射。设空间 B' 是仿紧的。若映射 \(f_0\) 和 \(f_1\) 同伦，则主 B'-纤维丛 \(f_0^*E\) 和 \(f_1^*E\) 同构。

推论 3. --- 设 B 是一个仿紧拓扑空间，G 是一个拓扑群。若 B 与一个点同伦等价，则每个以 B 为底的主 G-丛都是可平凡化的。

注记。 --- 这些结果的另一种证明方式，是验证命题 2 及其推论中构造的纤维空间同构实际上都是主纤维空间同构。

